\documentclass[nonflat,smallfont
% ,draftmode
]{langsci/langscibook} 
\renewcommand{\lsSeries}{guidelines} 
\usepackage{./langsci-tbls}
\usepackage{tabularx,multirow}
\usepackage{langsci-optional}


\renewcommand{\lsSeriesHistory}{~}
\renewcommand{\lsID}{opendata}
\renewcommand{\lsAdvertisement}{}
\renewcommand{\lsImpressumCitationText}{
  \onlyAuthor
  \renewcommand{\newlineCover}{}
  \renewcommand{\newlineSpine}{}
  Sebastian Nordhoff. %
  {\lsYear}. %
  \textit{Language Science Press business model}. % 
  Berlin: Language Science Press.
}
\BackBody{This book describes the business model of Language Science Press as devised in 2015. It gives some theoretical background and the text of the actual model itself. Each section of the model is then evaluated based on the knowledge of 2018, and possible alternatives which could have been adopted are discussed. An appendix with projected figures and actual figures for costs and revenue completes the book.

This book is complemented by the Open Access cookbook, open business data, and a spreadsheet for drafting and calculating own business models.
}


\title{Language Science Press business model}
\author{Sebastian Nordhoff\newlineCover Language Science Press}

\renewcommand{\lsISBNdigital}{978-3-96110-092-7}
\renewcommand{\lsISBNsoftcover}{978-1721068708}
\BookDOI{10.5281/zenodo.1286972} 

% %%%%%%%%%%%%%%%%%%%%%%%%%%%%%%%%%%%%%%%%%%%%%%%%%%%%%%%%%%%%%%%%%%%%%
%
%    Color definitions:
%
%%%%%%%%%%%%%%%%%%%%%%%%%%%%%%%%%%%%%%%%%%%%%%%%%%%%%%%%%%%%%%%%%%%%%

\definecolor{lsLightBlue}{cmyk}{0.6,0.05,0.05,0}
\definecolor{lsMidBlue}{cmyk}{0.75,0.15,0,0}
\definecolor{lsMidDarkBlue}{cmyk}{0.9,0.4,0.05,0}
\definecolor{lsDarkBlue}{cmyk}{0.9,0.5,0.15,0.3}
\definecolor{lsNightBlue}{cmyk}{1,0.47,0.22,0.68}
\definecolor{lsYellow}{cmyk}{0,0.25,1,0}
\definecolor{lsLightOrange}{cmyk}{0,0.50,1,0}
\definecolor{lsMidOrange}{cmyk}{0,0.64,1,0}
\definecolor{lsDarkOrange}{cmyk}{0,0.78,1,0}
\definecolor{lsRed}{cmyk}{0.05,1,0.8,0}
\definecolor{lsLightWine}{cmyk}{0.3,1,0.6,0}
\definecolor{lsMidWine}{cmyk}{0.54,1,0.65,0.1}
\definecolor{lsDarkWine}{cmyk}{0.58,1,0.70,0.35}
\definecolor{lsSoftGreen}{cmyk}{0.32,0.02,0.72,0}
\definecolor{lsLightGreen}{cmyk}{0.4,0,1,0}
\definecolor{lsMidGreen}{cmyk}{0.55,0,0.9,0.1}
\definecolor{lsRichGreen}{cmyk}{0.6,0,0.9,0.35}
\definecolor{lsDarkGreenOne}{cmyk}{0.85,0.02,0.95,0.38}
\definecolor{lsDarkGreenTwo}{cmyk}{0.85,0.05,1,0.5}
\definecolor{lsNightGreen}{cmyk}{0.88,0.15,1,0.66}
\definecolor{lsLightGray}{cmyk}{0,0,0,0.17}
\definecolor{lsGuidelinesGray}{cmyk}{0,0.04,0,0.45}
 
\definecolor{langscicol1}{cmyk}{0.6,0.05,0.05,0}
\definecolor{langscicol2}{cmyk}{0.75,0.15,0,0}
\definecolor{langscicol3}{cmyk}{0.9,0.4,0.05,0}
\definecolor{langscicol4}{cmyk}{0.9,0.5,0.15,0.3}
\definecolor{langscicol5}{cmyk}{1,0.47,0.22,0.68}
\definecolor{langscicol6}{cmyk}{0,0.25,1,0}
\definecolor{langscicol7}{cmyk}{0,0.50,1,0}
\definecolor{langscicol8}{cmyk}{0,0.64,1,0}
\definecolor{langscicol9}{cmyk}{0,0.78,1,0}
\definecolor{langscicol10}{cmyk}{0.05,1,0.8,0}
\definecolor{langscicol11}{cmyk}{0.3,1,0.6,0}
\definecolor{langscicol12}{cmyk}{0.54,1,0.65,0.1}
\definecolor{langscicol13}{cmyk}{0.58,1,0.70,0.35}
\definecolor{langscicol14}{cmyk}{0.32,0.02,0.72,0}
\definecolor{langscicol15}{cmyk}{0.4,0,1,0}
\definecolor{langscicol16}{cmyk}{0.55,0,0.9,0.1}
\definecolor{langscicol17}{cmyk}{0.6,0,0.9,0.35}
\definecolor{langscicol18}{cmyk}{0.85,0.02,0.95,0.38}
\definecolor{langscicol19}{cmyk}{0.85,0.05,1,0.5}
\definecolor{langscicol20}{cmyk}{0.88,0.15,1,0.66}


\newcommand{\lsptable}[2]{
\resizebox{#1}{!}{
\begin{tabularx}{\textwidth}{XXXXXXXXXXXXXXXXXXXX}
 \cellcolor{langscicol1}&\cellcolor{langscicol2}&\cellcolor{langscicol3}&\cellcolor{langscicol4}&\cellcolor{langscicol5}&\cellcolor{langscicol6}&\cellcolor{langscicol7}&\cellcolor{langscicol8}&\cellcolor{langscicol9}&\cellcolor{langscicol10}&\cellcolor{langscicol11}&\cellcolor{langscicol12}&\cellcolor{langscicol13}&\cellcolor{langscicol14}&\cellcolor{langscicol15}&\cellcolor{langscicol16}&\cellcolor{langscicol17}&\cellcolor{langscicol18}&\cellcolor{langscicol19}&\cellcolor{langscicol20}
 \rule{0pt}{#2}
\end{tabularx}
}
}

\newcommand{\background}[1]{ 
  \vspace{5mm}
  \renewcommand{\tblslinecolour}{lsDarkBlue}
  \tblssy[red]{explore2}{Background}{\vspace*{-5mm}#1}
}
\newcommand{\langscisolution}[1]{
  \renewcommand{\tblslinecolour}{lsLightBlue}
  \tblssy{langsci}{LangSci solution}{\vspace*{-5mm}#1}
}
\newcommand{\evaluation}[1]{
  \renewcommand{\tblslinecolour}{lsLightOrange}
  \tblssy{receipt}{Evaluation}{\vspace*{-5mm}#1}
}
\newcommand{\othersolutions}[1]{
  \renewcommand{\tblslinecolour}{lsDarkGreenOne}
  \tblssy{more}{Other solutions}{\vspace*{-5mm}#1}
}

\renewcommand{\tblssy}[4][black!12]{% Use like this \tblsfr[colour]{book|bulb|glass|law}{title}{text}, where [colour] is optional
  \renewcommand{\langscisymbol}{#2}\renewcommand{\tblsboxcolor}{#1}
  \begin{mdframed}[style=yellowexercise,frametitle={#3}]
    #4
  \end{mdframed}
}


\newcommand{\biberror}[1]{{\color{red}#1}}
\begin{document}
\maketitle
\tableofcontents
\mainmatter
\chapter{Introduction}
This document details the business model of an open access publisher, using Language Science Press as an example. It covers the
business case,
the market analysis, 
the internal organisation, 
and financing. 
For each domain, some general background is given, signalled by {\color{yellow}\includegraphics[height=1em]{langsci/graphics/explore.pdf}}.
Based on this, the solution chosen by Language Science Press is highlighted ({\color{lsLightBlue}\includegraphics[height=1em]{langsci/graphics/tbls-langsci.pdf}}).
That solution is evauluated ({\color{lsLightOrange}\includegraphics[height=1em]{langsci/graphics/tbls-receipt.pdf}})
and other possible solutions are discussed ({\color{lsDarkGreenOne}\includegraphics[height=1em]{langsci/graphics/tbls-more.pdf}}). 

The model presented here is the model as it was developed in 2015. At the time of writing (2018), the model has evolved. These changes will be listed in the evaluation sections. 

The original business model was written in German. The text here is a translation, with slight adaptions for readability or compliance with the sectioning. Some passages which were to heavily tied to the internal workings of Freie Universität have been removed, while other passages have been rendered more explicit for an audience lacking familiarity with the academic, administrative and legal landscape in Germany. 



\chapter{Summary} 
\background{ 
A business model should have a concise summary, where you give an overview of your goal, motivation, and methods
}
\langscisolution{ 
Language Science Press publishes high quality Open Access monographs and edited volumes in the domain of linguistics. It is directed by Stefan Müller (Freie Universität Berlin) and Martin Haspelmath (Max Planck Institute for the Science of Human History), who are supported by an international advisory board of renowned scholars.

The project was born from the insight that community-based publishing of scientific literature will be more cost-efficient than publishing by commercial publishers. New technologies mean and the collaborative work methods enabled by Internet mean that many of the traditional tasks provided by publishers are no longer needed or can be provided at a fraction of the cost. 

Furthermore, the inherent interest of maximizing profit margins is often at odds with the interest of science. This is first and foremost the case of business models which rely on monetizing access to information via paywalls. These models have to make sure that the information they provide remains a scarce resource. A publishing platform run by scientists themselves, which has as its stated goal the furthering of access to knowledge, rather than the restriction of access to knowledge, represents a valuable asset for society at large. 

Language Science Press has received seed funding by the DFG (2014/06 – 2016/05). Following up on this initial funding, the project is set to continue sustainable operations, accessing various revenue streams (individual membership, institutional membership, print margins, donations). 

The principles of Language Science Press are: 
\begin{enumerate}
\sloppy
 \item \textbf{Open Access}: all books are available without a fee, worldwide. Printed copies can be ordered for low prices via print-on-demand service providers. Additionally, all other aspects of the publication process should also be as open as possible (Open Source Software, Open Business Data, optional Open Peer Review)
 \item \textbf{Quality}: Peer review is mandatory. Language Science Press positions itself as a premium publisher with the corresponding requirements for manuscripts. There is a fair chance that a manuscript will be rejected, and authors should know this. 
 \item \textbf{Specialization}: All books appear in specialised series, which assess quality of the manuscripts according to the standards of the subdiscipline. There are no manuscripts out-of-series since quality could not be guaranteed in those cases. 
 \item \textbf{Decentral organisation}: Series editors operate worldwide and in an autonomous fashion. They are responsible for acquiring new manuscripts and assuring their quality. For this, they select an editorial board. Language Science Press starts with 5 series. New communities can form series of their own. They have to submit a detailed proposal, which is then evaluated by the advisory board. 
 \item \textbf{Community}: The scientific community is directly involved in the publishing process. Next to authoring or reviewing books, community members can also take over roles in proofreading, typesetting, illustrating, marketing via crowdsourcing. 
 \item \textbf{Automation}: The book production process will be highly automated, making use of templates and well-defined work flows and tool chains as far as possible in order to reduce costs.  
 \item \textbf{Lean Publishing}: As a new project, Language Science Press has the advantage that it can concentrate on Open Access publishing and does not have to deal with legacy domains such as: warehousing, royalties, paywalling, intellectual property rights management, book stands. 
\end{enumerate}

The main service provided by Language Science Press is the coordination of the publishing process and establishing a brand for high quality OA books in the book market. Further services include acquisition of authors, author support, typesetting, indexing, design, distribution, dissemination, archiving, enrichment. 
}
\evaluation{In hindsight, the principles we now put forward as central are: 1) Openness, 2) Community, 3) No frills.}
\othersolutions{The summary is obviously tied to the project, so there are as many summaries as there are project, possibly more than one per project. 
}


\chapter{Business case}
\section{Starting point}
\background{
Detail the situation found at the beginning of the project idea. 

}
\langscisolution{
In August 2012, the OALI initiative (Open Access in Linguistics) was set up by Stefan Müller and other linguists working at the Freie Universität Berlin. As a first step, the initiative contacted many reputable linguists from all over the world and invited them to join as supporters, which was met with great overall willingness. Together with Martin Haspelmath from the MPI-SHH, Stefan Müller submitted a grant proposal for te call  ``Open Access Monographs in the Humanities''. 
The DFG granted the project and provided 575.000€ for the creation of a business model and the setup of the initial platform. 
The project started in June 2014. The project team consisted of two programmers, one sysadmin, one coordinator and an economist. 

Language Science Press is organised in a decentral way. Editorial boards are reponsible for the individual series, which each cover a specific subdiscipline in the field of linguistics.  They are expected to deliver manuscripts which are in principle ready for typesetting. 

In order to lessen the burden of authors, and in order to assure a faster dissemination of scientific findings, we aim at a speedy publication, meaning a maximum period of 9 months between the initial submission and the publication online and in print. Given that the groundwork was already done before the official project start in 2014, the first volumes could already be published in 2014 itself. 

The grant proposal projected 5 series with 4 books each for 2014 (=20 total books); 7 series with 5 books for 2015 (=35), and 9 series with 6 books in 2016 (=54 books). 
 
As of 2015, Language Science Press is based at the Arbeitsgruppe für Deutsche Grammatik (led by Stefan Müller) at the Institut für Deutsche und Niederländische Philologie. Hosting is done by the Centrum für Digitale Systeme at the Freie Universität. This cooperation has proven successfull and should continue beyond the projected end of the grant
}

\section{Vision, mission, goals, indicators}
There is a somewhat structured protocol for starting an enterprise. It consists of putting your ideas on paper and crafting a) a long term vision, which you cannot achieve right away, but which will guide you;
b) a mission, which is ambitious, but which can be reached; 
c) the goals, which give a clear indication when your mission is accomplished; and
d) indicators which measure the progress you make towards your goals. 
This protocol is not limited to the corporate sector: any kind of polity could use it, e.g. trade unions, NGOs or political parties, and in fact, they do. 

\subsection{Vision}
\background{The vision should be long-term and should be concise. It is generally not necessarily within immediate reach. ``End world hunger'' would be a good example for a vision for instance.}
\langscisolution{
All linguistics content is available to everybody without a fee in a sustainable structure. Membership of a university or a particular country must not be a precondition for accessing the content.}
\evaluation{This vision is still current as of 2018.}
\othersolutions{One could have opted for a broader vision (``all scientific content'') or a narrower vision (e.g. drop the second sententce of the vision). As of 2018, we feel that the vision is still sufficiently far away that it was not too easy a vision, and at the same time, it is not completely unrealistic that this vision could be reached one day, so we are happy.}
 
\newpage  
\subsection{Mission}
\background{The mission describes the raison d'être of your enterprise. Why are you here anyway? It should be possible that you accomplish your mission, even if it is bold.}
\langscisolution{ 
We will provide a platform for the distribution of world-class research in linguistics where the scientific community can distribtute their research results as open access books in a maximally autonomous fashion.
}
\evaluation{
As compared to the vision, the mission does not cover articles or other non-book content. Furthermore, the mission only covers the organisational setup of a platform; it does not cover lobbying for policy changes or the like. This does not entail that those would not be laudable goals. But they do not form part of the core mission of Language Science Press. Other actors might have a vision which includes e.g. lobbying for policy reform. This is not necessarily in conflict with our mission.

We are regularly asked whether we will host journals or provide other services. The stated mission helps us to remain on track and not lose our focus.}
\othersolutions{One could have opted for a less technical mission and have a more prominent focus on policy change; one could have had a similar mission for article or for research data; one could have focussed more on Citizen Science. These other missions would not be in conflict with our mission. } 

\subsection{Goals}
\background{In order to not get lost on the way while working on realising your mission, it is important to sketch the itinerary of your journey and the goals you want to reach. The (shortish) mission can be translated into a number of more concrete goals. Where the mission is bold and vague, the goals are achievable, specific, measurable, and timed. There should be a way of evaluating whether you have made any progress towards your goals.}
\langscisolution{
We have fixed the following goals
\begin{itemize}
\item High customer satisfaction (Readers, authors, editors, funders)
\item Creation of a strong brand owned by scientists which stands for quality and innovation
\item Coverage of all subdomains of linguistics
\item Continuous output
\item Solid and sustainable setup
\end{itemize} 
}
\evaluation{
The 5 goals are achievable. They are specific in the sense that you know whether you have reached them or not. They are not particularly well suited for being measured with the tools we have at hand. For instance,  we lack the means do any surveys for customer satsifaction or to measure the strength of the brand. ``Solid and sustainable setup'' can be somehow evaluated, but it would be difficult to apply quantitative measures. The only goal which really meets the criteria is ``continuous output'', where one can easily count the number of publications, and measure the intervals between them. This being said, the absence of a quantitative evaluation does not really seem to have been fatal to the usefulness of our goals. 
}
\othersolutions{
Other possible goals would be market penetration, just representation of different demographic groups, attraction of high quality submissions. These were not among the stated goals, but if those goals happened to be achieved as well, that would not  be problematic.
} 


\subsection{Indicators}
\background{How is progress being evaluated? For every goal, find a dimension where you can measure your progress towards that goal. That dimension will be an indicator. Typical dimension would be ``items sold'', ``revenue'' or ''time needed to perform X''.}
\langscisolution{
The following list gives our indicators, with the goal they relate to in parentheses.
\begin{enumerate}
\sloppy
\item Authors and readers are happy (Customer satisfaction) 
\item Stable series with regular new publications (Customer satisfaction: editors)
\item Acquisition of funds  (Customer satisfaction funders)
\item Download figures  (Customer satisfaction, strong brand) 
\item Submission figures (strong brand)
\item Copying of LangSci processes and models in other disciplines (strong brand)
\item Rising number of non-overlapping series (strong brand, coverage)
\item Rising number of publications (Output)
\item Language Science Press has full legal capacity and clearly defined headquarters (solid setup)
\item Stable staff with adequate payment (solid setup)
\item Stable and long-term relations with service providers and clients (solid setup)
\end{enumerate}
 
}
\evaluation{
The indicators given above have for the most part been useful. More details are given below.
\begin{enumerate}
\item difficult to monitor, but done qualitatively
\item output per series is monitored and ailing series receive support
\item one project proposal was submitted since this model was written. It was granted. In that sense, this indicates that is being monitored, although the one data point itself does not allow for detailed interpretation. 
\item Continuously monitored and evaluated. Outliers are investigated (It turns out that textbooks are above average, and non-English books are below average, which is both easy to explain). 
\item Continuously monitored as for quantity and regularity. 
\item No information available. Most likely due to no one copying the model. Not particularly useful as an indicator.
\item Monitored. Useful to a certain extent, up to saturation. 
\item Monitored. Useful. 
\item Monitored. Legal form is completed, but headquarters are still under development. Although this indicator is qualitative, it is very useful. 
\item Monitored. No fluctuations. Staff seems to be very happy with the employer, but this is not quantified. Useful indicator. 
\item We have a number of returning authors and many returning proofreaders. Our relation to the BoD print-on-demand service provider was intensified. Useful indicator.
\end{enumerate}
}
\othersolutions{ 
Other indicators which one could have chosen, but which were not part of our set would be:
website visitors, print sales, followers on social media, Altmetrics, Impact Factor. The former three are actually being monitored, while the latter two are not. This is mainly due to linguistics being a discipline where bibliometrics play a lot less important a role than in other disciplines.
}

 
\clearpage
\section{Customers and value propositions}

\background{Describe who would be interested in your project (the customers) and why (the value the project represents for the people interested). This leads to the topic of how the people interested can be made to participate in the financing of a project which is valuable to them. 
}
\langscisolution{
The most obvious way of recouping the costs of book production is selling the books (``reader pays''). This option is not available for OA publishing, at least not in a straightforward way. The digital versions have to be available for free by definition, and the number of printed copies sold should not be overestimated. 

This implies that we have to separate the roles of ``reader'' and ``funder''. In terms of financing, the question to ask is: who has an advantage from OA publications? Who will derive a benefit if content is available for free online? In business lingo, who are the ``customers'', and what is the ``product'' being offered? 

[A short digression is in order here. In order to speak about financing models, one has to agree on a vocabulary. As it happens, there is already a vocabulary in place for discussing the concepts, actors, and processes which arise when setting up shop. This is the vocabulary employed by economists. In their terms, a ``customer'' is someone who derives a ``benefit'' from the existence of a particular ``product'' and is willing to transfer resources to the ``producer'' so that the customer can make use of the product to enjoy the benefit. Many linguists (and other scientists) will object to them being seen as ``customers''. This objection is understandable: the relation a researcher has with a book is very different from the relation the same researcher has with their supermarket or their phone company. This being said, it is necessary, in science and elsewhere, to clearly define the concepts, and to have labels to unambiguously refer to the concepts defined. In the context of this document, these terms  should be taken to be valid within the definitions employed here, and they should not be taken to contain any connotations that ``customer'', ``producer'', ``product'' and ``benefit'' might have in everday life, or in the economic organisation of society at large.]

The books being published by Language Science Press come from authors of countries all over the world, and are read internationally as well. This means that the business model must reflect this international character of the operations. A funding model restricted exclusively to, say, Germany, will be unfair, and it will be hard to convince the residents of only one country to fund operations which profit the whole world. This international setup has to be taken into account not only for financing, but also for distribution and marketing. 
 
When we analyse the benefits OA publishing provides, we find that it is not only the readers who profit. Also the authors derive a benefit (more prestige, faster publication, greater reach). Furthermore, libraries also do benefit via a larger array of books they can provide to their patrons. The state also benefits (more cost-effective publishing; no dependency on monopolists), as does society at large (better and faster access to scientific findings.) 

Just as it would be unfair to have one country shoulder the operating costs alone, burdening only one group with the costs (e.g. via ``author pays'' models) is questionable as well. So the question is: how can the operational costs be shared among the stakeholders in such a way that everybody contributes, and that no one is charged beyond their capacities? And how can that be done without undermining the basic tenets of Open Access?

Language Science Press has opted for a broad basis for financing its operations. The diversification of revenue streams means not only that each of the various stakeholders participates, but it also reduces the risk of funding gaps as compared to one large funder. If one revenue stream breaks down or does not contribute the expected amount, this can then potentially be compensated by the other streams. 

The analysis of ``clients'' above shows that the ``product'' provided by Language Science Press is complex and can access a variety of ``markets'' with different target groups. 
 
The business model presented here is based on five pillars, which are supported by the different audiences. For every audience, different services and extras are offered. The basic assumption is that each group of stakeholders will have have both egoistic motivations (purchase a book for reading it yourself) as well as altruistic motivations (create societal benefit, ensure free access to the knowledge of humanity, ...).  We acknowledge both kinds of motivation as valid. 

For long-term sustainability, Language Science Press requires support from institutions, which provide funding on a regular basis via memberships or BPCs, but also from individuals. The latter provide less finanical support and more support in kind, but they are important for reach out and dissemination. Only when a broad base of linguists is aware of Language Science Press will the prestige go up and the project can transform into a publisher in the strict sense. 

A short version of the individual strategies for the five ``pillars'' can be found in the appendix. The particular target groups, the service packages (products, services), marketing and projected revenues are listed there. More elaborate strategic discussions can be found in the cookbook accompanying this business model. Real world data are also provided for inspection together with thes documents. 

As far as strategies go, it is important to note that OA monographs represent a field where little prior knowledge and experience is available. Funding OA via an array of different and voluntary revenue streams is certainly an innovative idea, but it lacks empirical underpinning. Particularly in what regards the amount of revenue effectuated, only broad guesses are possible. There are some studies around, but they are few and far between, and more often than not cover other fields of science (e.g. natural sciences) or ``products'' (e.g. articles). Therefore, the concept presented by Language Science Press should be tested for two years. It should then be evaluated, and adjustments according to the outcomes of the evaluation. 
}
\evaluation{
The analysis of the stakeholders still holds true. In the model we claimed, however, that we would go for a diversified funding model. As a matter of fact, of the five ``pillars'' identified, only institutional membership is reliable, steady and able to cover an important share of the costs. The other four pillars provide a little extra money, but they would not be able to sustain operations if institutional membership did not exist. As a result, we can say that the ``broad basis'' is actually not so broad, and that the largest burden is being carried by libraries. (Note that under the traditional model, it is also libraries which pay the largest share).
}

\othersolutions{
One could have opted for a narrower approach right away, e.g. only go for BPCs or only for institutional memberships. One could also have tried harder to get a more balanced distribution, instead of simply settling on institutional memberships as the most promising pillar. 
}


\section{Technology and domain knowledge}

\background{It is important to take stock of the domain knowledge (linguistics and publishing in our case) and the technological skills and hardware or software available when conceiving the project and required when starting. For instance, it could be the case that the enterprise is heavily skewed towards one particular subfield. In that case, finding partners from other areas would be good idea to have a broader base. Or it might be the case that technical typesetting skills are available but no skills regarding webserver administration. In that case, appropriate staff have to be hired or service providers have to be identified. Finally, large monitors with high resolution with be a technical requirement which cannot be taken for granted, and specialized software might have to be bought.}
\langscisolution{
The technological skills required for this project are almost exclusively in the domain of IT. The development and maintenance of the platform and of the publication process do not require infrastructure beyond off-the-shelf desktop computers. Since books take up very little space on the hard disks, and distribution does not make use of complex algorithms demanding powerful processors, either, a small server is sufficient. 
On the software side, we can mention a web platform for creation and distribution of books. The existing platform OMP fulfills these requirements and is available without cost as open source software.

The Language Science Press team has a large breadth of experience with scientific publication. Both press directors, Haspemath and Müllers, are prolific in their scientific output. Stefan Müller has been editing the proceedings of the HPSG conference since 2003 (12 volumes in total as of 2005). Martin Haspelmath has edited 7 volumes since 1995 and has pioneered the domain of ``expert consortia'' in linguistic publishing with projects such as WALS, WOLD, APICS, ValPal, with project teams numbering beyond 70 participants. 

Sebastian Nordhoff as the project coordinator founded the field of Electronic Grammaticography and has wide experience in electronic publishing and digital humanities as a whole. This concerns both data publication as well as publication of monographs and edited volumes. He has devised the platform ``Enhanced Electronic Grammars'' for de Gruyter, specified the back end and converted 20 large and complex grammars for ingestion.
}
\evaluation{The initial description proved accurate. Skills which were missing in the beginning were vector drawings (e.g. maps) and screencast production. Also, the social media skills could have received more attention in the beginning.  A need for more formal project management skills was identified and Nordhoff received training in this domain.}
\othersolutions{The approach followed by Language Science Press was very much to develop the necessary skills as the need for them arose (vector maps, video editing, chart generation). One could instead have opted for a much more formalized approach with key competences and development strategies for the staff. Since there is only one coordinator and one student assistant, this was judged to be too much overhead.}


\chapter{Market}

\section{Market situation}

\background{It is important to have a clear idea of the market where your enterprise is going to operate. This is true even if yours is a non-profit enterprise, since some of the competitors might be for-profit. }
\langscisolution{
The market for scientific publishing has been very dynamic in the recent years and continues to be in a state of disruption. On the one hand, we witness the change from print to digital, on the other and, the traditional subscription-based model is being replaced by novel concepts (e.g. ``author-pays''). 

There are several ways to split the market: 
\begin{itemize}
\sloppy
 \item according to medium: market for scientific articles; market for monographs; market for research data
 \item according to subject: natural sciences vs. humanities; or, more detailed, according to discipline
 \item according to distribution model: closed access vs. open access.
\end{itemize}
 
The different submarkets listed above show only partial overlap. Some disciplines have very speficic traditions and communities of practice when it comes to publishing. For instance, the OA market for articles in the natural sciences can be called mature. OA journals and platforms are well-established and integrated into the practices there. In the humanities, this area is less populated and there is more room for development.

In the smaller domain of OA monographs, there are only very few projects across all disciplines. This is probably due to the trend towards ever shorter publications in the natural sciences, to the detriment of books. Monographs have been declared dead several times over the course of the last years, and the OA movement has focussed nearly exclusively on journals and repositories in the past. However, it has turned out that the commnunities of practice in the humanities will not switch that easily to the publication of articles as the sole medium of knowledge interchange. Longer, more encompassing, works will remain indispensable in the humanities. With the backdrop of the change towards OA sketched above, a change in this domain is unavoidable as well. Language Science Press can be a pilot here, and have a significant first-mover advantage. 
}
\evaluation{
We still stand by this analysis. Events which have an influence here but had not been considered are the Lingua/Glossa-transition and the rise of the Open Library of Humanities; the replication crisis and calls for Open Science; the DEAL negotiations in Germany.  
}
\othersolutions{
What is completely missing in the market analysis is the difference between modern and legacy models of book productions. Traditional publishers are very much tied to a first proofs/final proofs model, whereas publishers born digital can accommodate a much more fluid model, and openness means that no provisions must be taken to shield the content, allowing open peer review and community proofreading on platforms like PaperHive or docLoop.

Furthermore, this market analysis does not at all take into account possibilities of monetization or analyse competitors, their strenghts and weaknesses.%TODO
This is basically due to LangSci trying to provide a service for the community by the community, a service which was lacking. The purpose is not at all to dislodge another actor, but simply to make sure that academics have a way of making their research available which is according to their needs. Furthermore, the driving force is not the desire to make money, which might be why this aspect was simply overlooked.
}
 
\section{Target audience and customers}

\background{It is important to know who has a need for the product or service you are providing. Try to be as specific as possible, and try to get an idea how large the relevant group(s) would be. Then, try to find a way how these groups can contribute towards the running of your enterprise, either financially or by providing services in kind.
}

\langscisolution{
For Language Science Press, we can identify four target groups, with different roles. Language Science Press will provide customized services and extras. At the same time, we assume that the groups will also have idealistic motivations and want to instrically support the idea of Open Access. Therefore, material ``benefits'' (e.g. in the form of physical books) are only one part of the ``customers''{}' motivation. 

{\small
\noindent
\begin{tabularx}{\textwidth}{l>{\raggedright}Qp{42mm}}
\lsptoprule
\textbf{Group} &  \textbf{Customer relation} & \textbf{Services \& extras}\\
\midrule
authors &  Optional BPCs & Professional publishing process; OA publication\\
\midrule
libraries & purchase of printed copy & High quality printed copies as a complement to free online pdf\\
\midrule
\multirow{2}{33mm}{institutions (universities, resesarch centres, funders, learned societies)} & institutional membership& extras and special offers\\
\tablevspace
\cmidrule{2-3}
& donations & Promotion of the idea of OA\\
\midrule
\multirow{3}{33mm}{individuals (researchers, interested citizens, activists)} & purchase of printed copy& High quality printed copies as a complement to free online pdf\\
% \tablevspace
\cmidrule{2-3}
& individual membership & extras and special offers\\
% \tablevspace
\cmidrule{2-3}
& donations & promotion of the idea of OA\\
\lspbottomrule
\end{tabularx}
}


Many OA projects list authors as the most important source of revenue. Language Science Press also intends to raise some money from the group of authors since this group profits the most from the increased reach. But the funding landscape is such that author fees for books are all but non-existent in the funding programs of the relevant funding bodies in Germany and elsewhere. Therefore, any authors fees must be purely optional. Otherwise, there would be the risk that certain authors would be preferred not because of the scientific quality of their books, but because of their willingness to pay.  Platinum OA, as mentioned in the initial grant proposal for the DFG remains a fundamental requirement. There are two options to set this up: either have generally no fees unless an author wants to pay (``pay-what-you-want''), or have general fees with a no-questions-asked-waiver for anybody who feels they cannot afford the fees. Depending on the legal context, one or the other option might prove more promising.

Libraries are mainly relevant as buyers of printed copies. The library budgets should not be overstretched here, hence the printed copies will be provided at a low price. 

As a general principle, funding from a plurality of different institutions (universities, research centres, funders, learned societies, ...) is preferable. If the costs are shared between many funders, chances are that we will see similar projects in the future. This would profit all sides.

Scientists, interested laypeople and OA activsits are welcome as indivdual supporters. Some members of this group have already asked how they could support Language Science Press financially, with small but regular contributions. Unrelated to their being affiliated to a research institutions, people can find the idea of OA attractive and are willing to spend some of their money towards this goal. Both memberships as well as one-off donations can thus play a part in the financial success. More important than the financial contribution, however, is the time and energy supportes devote to Language Science Press in writing, reviewing, and proofreading.

We foresee to provide extras to supporters, which will complement and extend the free content. This could be access to early versions, free copies, background information from behind the scenes etc. 
}
\evaluation{
Our description of the target audiences lacks numbers. It is completely unclear how many authors, readers and libraries would be in the relevant groups. 

Our ideas regarding extras and special offers for libraries and authors were completely mistaken. We initially had thought about some kind of benefits for libraries, e.g. lower author fees for participating institutions, or free copies for supporters or the like. It turns out that there is no demand for these kinds of extras. Authors, readers,  and libraries both want to support us on idealistic grounds, and having some kind of extra or not will not make a difference there. In the worst case, it will be met with bewilderment about what the ulterior motives behind giving away the extra would be. As for libraries, it might also put them into ``spreadsheet mode'', where they will start calculating whether becoming a member at cost $X$ a year is actually worthwhile when they have only $n$ authors publishing with Language Science Press a year. The major argument for supporting Language Science Press, however, is not the desire to save a couple of hundred Euros in author fees, but to bring about a change in the publishing landscape. Focussing on the savings of Euros will distract from the main message. 

The capacity or willingness of individuals to provide financial support has been vastly overerstimated. Researchers are more than willing to provide effort in kind (e.g. reviewing or community proofreading), but financially, only very low sums could be raised (see below). 

As a whole, the plurality of different institutions willing to fund has been overestimated. 

The options ``pay-what-you-want'' and ``no-questions-asked-waiver'' have not been investigated in detail since there were very few books with processing charges in the first place. 


}
\othersolutions{Instead of having a broad approach and trying many different target audiences, one could have concentrated on only one group with higher precision in the messaging and reduced overhead. Or, one could have tried to establish a product which would be interesting to all groups alike in order to simplify processes. Finally, one could also have a book-centered approach, where people can pledge not for the platform as a whole, but for individual books. As soon as the nominal BPCs for a book a reached via that pledging, production will start, similar to crowdfunding for music albums or computer games. This is, however, not practical for scientific books. First, you do not want to have the most popular books only. Especially for science, it is important that books like ``A grammar and dictionary of Chakali'' are also produced, even if they lack wide popular appeal. Second, bookkeeping of all the pledges will be too much overhead. Third, the main costs are infrastructural costs. The costs per title are not so high once the infrastructure is set up. This means that we have to fund the infrastructure, not the extra cost per title. Fourth, we employ staff. We cannot tell them that they will not get salary until the next title reaches its pledging target and that until then they are on leave.}

 
\section{Competition }
\background{Who are the major players, what are their strategies, what are the latest developments, what are the demands of the ``clients''?}
\langscisolution{
In contradistinction to journal articles, there are as of 2015 no established methods of publishing linguistic books in Open Access. Commercial publishers (for linguistics: de Gruyter, John Benjamins, Oxford University Press, and Cambridge University Press) do have lukewarm options on offer, but do not have a well-defined program, nor do they really publicize those options. The Open Access options by these traditional publishers are only a variation of their more classical offers, and the costs of more than 10,000 EUR is prohibitive for most authors. A dedicated branding strategy or a specialization/division of labour between Closed Access and Open Access product lines cannot be discerned.

A community-based model with autonomous series, as embodied by Language Science Press, is currently non-existent in the book sector. This means that Language Science Press has the possibility to claim terrain here and acquire standing as a prestigious publishing platform for OA books.  The brand (implying prestige conferred by the publisher) remains up to this day one of the main criteria which determine authors' choices of a publishing venue since their future career will heavily depend on this form of symbolic capital. 

In the technical domain, the competition is caught in their legacy technology and lacks the flexibility to adapt to novel trends or requirements. For instance, an innovative electronic product devised by Nordhoff for a large German publisher had to be integrated into the existing content management system, which stripped the platform of its intended key features. 

The traditional publishers do have advantages emanating from their centuries-old experience in the print sector, but they are less flexible. As far as the integration of the scientific community goes, the lead Language Science Press has built up will be difficult to catch up by the commercial competition.  

Other Open Access projects evolve in the same domain (e.g. Open Library of Humanities, Edition Open Access, Topoi, Cleo), but the relation is amicable here, rather than competitive.
If these projects were to offer service in the long run which would overlap with services provided by Language Science Press, the resolution would involve collaboration rather than competition. As of 2015, these projects are not active in the field of books in linguistics. 
}
\evaluation{This analysis still stands as of 2018. ``De Gruyter Open'' has been discontinued in the meantime and rebranded as ``Sciendo''. }
\largerpage
\othersolutions{One could have broadened the analysis to include other scientific domains (e.g. History or Law), or other forms of publication (journals). As for the competitors, one could also have made a more detailed analysis of the mentioned publishers and their inferred strategies in the OA market.}
 
\section{Geographical location}

\background{Think of criteria which determine where your enterprise should be located. Proximity to resources, clients, workforce and relevant political actors are important criteria.}
\langscisolution{
Berlin is well suited as a base for an innovative publisher. We can cite PaperHive or ScienceOpen as other Berlin companies active in the area of Open Access. In the non-profit domain, we can list the Wikimedia Foundation, the Open Knowledge Foundation, and the Chaos Computer Club, which highlight the vibrant societal basis for projects involving open projects. 

The Freie Universität Berlin with its focus on the humanities is well suited for this project. Its slogan ``International Network University'' underscores the appeal for a publishing platform bringing together scientists from all over the world. 

The close connection between the ``Arbeitsgruppe Deutsche Grammatik'' at the ``Institut für Deutsche und Niederländische Philologie'' and the ``Centrum für Digitale Systeme'' is also an asset for the Freie Universität, especially since the CeDiS has the hosting of Open Journal System (OJS) and Open Monograph Press (OMP) as one of its core competences.
}
\evaluation{This analysis is kind of bogus since for personal reasons, a base other than Berlin would never have been considered. Still nice to have a writedown of the key advantages, though.}
\othersolutions{London or The Netherlands would have been other obvious choices. One could also have dropped that section altogether and position the enterprise as a ``virtual publisher'' with headquarters on the Internet and no specific geographical location.}
 

\section{Marketing and distribution}

\background{
Think about how your target audience will know about your product or service, and how it will be made available to them. 
}
\langscisolution{
The press directors have a very extensive and international network within the scientific community, which is evidenced by the high number of supporters. As of 2015, Language Science Press lists more than 500 supporters on its website, among which more than 160 professors. This network can be activated, and the number of new series proposals and book proposals underscores the demand for a true OA publisher in linguistics. 

The following output is projected:

\noindent
\begin{tabularx}{\textwidth}{Xrrr}
\lsptoprule
        & series &  books per series & books total   \\
\midrule         
Year 1 & 16      &  3                & 48              \\
Year 2 & 18      &  3                & 54              \\
Year 3 & 20      &  3                & 60              \\
Year 4 & 22      &  3                & 66              \\
Year 5 & 24      &  3                & 72              \\
\lspbottomrule
\end{tabularx}

Social Media and other electronic means of PR are used extensively for marketing and distribution and reflect the international setup of Language Science Press. Distribution and marketing are tailored for the different audiences (see Appendix) and will in part be covered by the communities themselves (book reviews, word of mouth, social media). This means that the activiation of the community is fundamental. This makes use of a cascading model: the LangSci core team takes care of finding active series editors, which in turn activate the members of their respective communities.  
}
\evaluation{
This section is basically paying lip service to the model in putting down some paragraphs about marketing. As a matter of fact, Language Science Press does not really do marketing. We believe that marketing is mainly needed to sell books people do not need in the first place. 

As for distribution, Internet makes it very easy.

What is missing from the model is how libraries will be approached. This was done in a two-tiered fashion. Supporters from linguistics were asked to get in touch with their library about Language Science Press. As a follow-up, Knowledge Unlatched then contacted those libraries a couple of days later. In that way, we attempt to make sure that we only contact the libraries which have a genuine interest and do not spam the rest.

The output numbers were overestimated. See the Appendix for a comparison of projected and acutal output. 
}
\othersolutions{ 
One could invest much more in marketing and try to promote the enterprise and the brand among the target audience, raise awareness, and generate higher traffic. This could be done with ads on the Internet, book stands, mailings, sponsorship or the like. 

If one chooses a reader-pays model, one would have to devise a way how fees for accessing the content will be collected, and how it can be made sure that only people entitled to access are actually included in the distribution.  Under an author-pays model, some workflow for the financial transaction has to be established, too, but this will be much simpler. 
}

\chapter{Organisation / Staff}

\section{Legal form}

\background{In the beginning, an scholar-owned publisher can operate without being a legal entity in its own right. The university normally can provide office space and servers, and the time needed to set up the operations is either included in the working time of the researchers, or they do it in their free time. 

Eventually, there will be a moment when some services will be bought from outside. This might be a designer, a web developer, PR, or print-on-demand services. In order to participate in such business dealings, the enterprise has to have a so-called ``legal form'' in order to enter into contracts, receive payments, pay taxes etc. The legal forms available depend on the jurisdiction, but generally they fall into a number of types:

\begin{enumerate}
 \item ``commercial section'' of a university
 \item Foundation 
 \item Association 
 \item Private company
 \item Registered company 
 \begin{itemize}
  \item charitable registered company
 \end{itemize}
 \item Public company 
\end{enumerate} 
}
\langscisolution{
Language Science Press shall have the legal form of ``Betrieb gewerblicher Art (BgA)'' ($\simeq$ ``commercial section of the university'') at the Freie Universität.
This reflects it close association with the scientists as well as the participation in a market where Language Science Press competes with commercial actors. 
}
\evaluation{
The administrative dealings with the unversity in order to enter into contracts with service providers, to send out invoices, and to pay takes proved a nightmare. For instance, it is completely impossible to use a company credit card when being part of a university. Deciding on the tax rate (0\%, 7\%, 19\%) took months of discussion, and it was unlear who would have the final say in these matters. 

As a result, \textit{LangSci Press gUG (haftungsbeschränkt)} was funded by Müller, Haspelmath and Nordhoff. This is a registered company with charitable goals (furthering of science and education). As such, it enjoys a number of fiscal privileges, but it is barred from  distributing profit to its owners. Nordhoff is the representative of the company and can now easily enter into contracts with service providers, get legal counsel for a fee, or send out invoices as required by the daily business. 
}
\othersolutions{
A foundation (\textit{Stiftung}) is highly regulated under German law and requires an important initial endowment. This capital was not available. 

An association (\textit{Verein}) is a more democratic legal form than a company. Associates have an Annual General Meeting, where they discuss the direction the association should take. This legal form was not flexible enough for our purposes and entails some unneeded administrative overhead. 

A private company (\textit{GbR}) is not listed in the  commercial register. A private company cannot have its charitable aims legally recognised in Germany, so we did not pursue this option. 

A public company (\textit{AG}) requires an initial investment of 50,000 EUR. This capital was not available. 
}
 

\section{Organisation}

\background{Based on the legal form chosen, describe the internal organisation of the enterprise. How are decisions made? Who is implicated? How are conflicts resolved?

What are the relation to external entities?}
\langscisolution{

The organisational structure after May 2016 shall continue the structure as of 2015 to a large extent.

Scientific oversight will remain with Stefan Müller and Martin Haspelmath. For the vetting of new series, there is the international Advisory Board (24 researchers worldwide). The individual series are led by one or more series editors each. All mentioned researchers do not receive remuneration from the project. 

Organisational and legal representation is given by Stefan Müller (5\% of his time).

The project team should have two coordinators (total 1.5 FTE), who are responsible for author support. An accountant (0.5 FTE) shall take care of finances, fundraising and implementation of the model; further support is given by a secretary (0.25 FTE) and a student assistant (0.5 FTE). 

Hosting will continue to be provided by Cedis

In order to provide its services, Language Science Press has to enter into legal agreements with persons external to the university: Authors sign an author agreement, where they license their work under a Creative Commons License; where they agree that Language Science Press distribute the work; and where they assure indemnity from liability to Language Science Press.

Print-on-demand service providers will be granted non-exclu\-sive rights. 

As for revenue, no long-term agreements will be entered into with institutional members and individual members. In order to reduce bureaucracy, institutional membership can be cancelled on a monthly basis. 

No further long-term agreements will be entered into. 

}
\evaluation{
The creation of the charitable company supersedes this section. The company has 3 associates and requires a 3/4 majority to take decision, which implies unanimity. If a fourth associate was accepted, unanimity would no longer be required.

The project now has only one coordinator (1.0 FTE), no accountant, no secretary, and a student assistant for 10h a week. As discussed below, the initial business plan had 50 books a year, while we now operate on 30 books a year.

The contracts with institutional members now have a 3-year term and cannot be cancelled. This reduces bureaucracy on all sides. 

The agreement with Knowledge Unlatched is also for three years, as is the contract with BoD.

A final solution for the hosting situation still has to be found.
}
\othersolutions{A higher output would entail a higher income, higher workload, and more staff. 

In the company, different shares of the associates (e.g 50:25:25) would have entailed different power relations.

One could have opted for exclusive Print-on-demand providers. 

One could have retired the Advisory Board after the initial phase. 

}
 
\section{Partners}

\background{Most enterprises can not achieve their mission on their own; they require interaction with outside entities, the partners. It is good practice to take stock of which other entities are essential or helpful for the achievement of the mission.}
\langscisolution{
The most important partners for Language Science Press are the series editors. Their work is not remunerated. Further cooperations exist with:
\begin{itemize}
 \item print-on-demand service providers: Book on Demand, Epubli, CreateSpace 
 \item Library of the Freie Universität 
 \item Various Open Access associations 
 \item Learned societies in linguistics
\end{itemize}
}
\evaluation{
What is completely missing in the description above is the community. We have realised that the greatest LangSci asset is actually our network of supporters, who help us with many tasks, ranging from community proofreading to PR or acquisition of institutional supporters. 

Other partnerships which are not listed above because they only developed in the course of the project are Knowledge Unlatched for dealing with institutions and PaperHive for online annotations and work flow. Another important partner is Overleaf for drafting, editing and revising manuscripts online in a collaborative fashion. All the listed partnerships would be very hard to replace if ever they were to fail.

Since the business model was conceived while LangSci was still based at the Freie University, that university is not specially listed; nevertheless, the Freie Universität remains an important partner for hosting and archiving even after the transition to Humboldt University.
}
\othersolutions{ 
Instead of relying on the community, one could have set up a network of freelancers who provide services in copyediting, typesetting or programming for a fee. In order to pay for that, one would have had to find additional sources of revenue. 

We also contacted the Open Library of Humanities about setting up a funding model together, but it turned out that their own operations were already very demanding and they would rather not take in additional processes from outside. 
}
 

\section{Team members}

\background{Detail who will be on the payroll. }
\langscisolution{
The organisational and legal project lead is taken over by Stefan Müller, in whose institute Language Science Press is based.

The scientific leader is Martin Haspelmath. While not based in Berlin, but in Leipzig/Jena, he is a key member of the Language Science press team.

Since March 2014, Sebastian Nordhoff has been the coordinator of Language Science Press. His focus lies on the improvement of workflows and on the usability of the website and other platforms. As the project progresses, Nordhoff will move on to support series editors and take over the management of Language Science Press at large.

A fourth person, to be employed, will take over coordination and support of the series and series editor not under Nordhoff's oversight. 

Accountancy and fundraising
[description of economist's experience redacted for reasons of  privacy law]

A secretary will take care of bookkeeping for donations, membership fees, and general correspondence.

Finally, the series editors are not part of the core team, but they are indispensable for the success of Language Science Press. A good deal of the classic tasks at a publishing house, such as acquisition, author support or marketing is carried out by them. 

}
\evaluation{The revenue model was radically simplified by relying almost exclusively on institutional memberships. The whole process of promotion, acquisition, invoicing and payment was outsourced to Knowledge Unlatched as a service provider, against a commission. This meant that the positions of the economist and the secretary could be dispensed with, and that some of Nordhoff's time could be shifted from liaising with institutions to book production. 

Nordhoff's position constitutes a single point of failure at this moment. In order to arrive at a more resilient structure, the position of the coordinator will be split in half among two persons, so that there is always a fallback available.
}
\othersolutions{
Instead of outsourcing the library liaison to Knowledge Unlatched, one could have done that in-house. This would have meant that an additional team member would have to be hired and be trained in the relevant legal and fiscal aspects, as well as in the communities of practice of librarians. 

One could also have tried to run the project with professors alone, i.e. without a dedicated coordinator. At about 200-300 emails per book, however, this would quickly surpass the time professors have at their disposal. 
}
 

\chapter{Funding}

\section{Budgeting}

\background{Your budget should include personnel and non-personnel costs. Some items will present fixed costs, while others depend on other parameters, e.g. the number of books.}
\langscisolution{
The major part of costs consists of costs for staff. The basis for calculating these costs is the German payscale TV-L with an annual preciation of 3\%

The operations of Language Science Press will grow. Every year, new series should join, and the network and the platform should become wider. Due to this quantitative increase, personnel costs will go up; at the same time more revenue will be generated. The following table gives the demands for staff:% and the costs: 

% \todo[inline]{costs missing}
\noindent
\begin{tabularx}{\textwidth}{Xrrrr}
\lsptoprule
         &  scientist  & secretary &   professor &   student assistant  \\
\midrule
Year 1 &  2,00 FTE                 &    0,25 FTE                       &   0,05 FTE    &   0,5 FTE                  \\
Year 2 &  2,00 FTE                 &    0,25 FTE                       &   0,05 FTE    &   0,5 FTE                  \\
Year 3 &  2,00 FTE                 &    0,50 FTE                        &   0,05 FTE    &   0,5 FTE                  \\
Year 4 &  2,25 FTE              &    0,50 FTE                        &   0,05 FTE    &   0,5 FTE                  \\
Year 5 &  2,25 FTE              &    0,50 FTE                        &   0,05 FTE    &   0,5 FTE                  \\
\lspbottomrule
\end{tabularx}

Non-personnel costs do not include taxes. 

Next to personnel and non-personnel costs, we also have to include calculatory costs resulting from being based at a university (e.g. office space or IT). This is modelled by taking over the costs from a research of the relevant posts on the Berlin market. 

Revenue also excludes taxes. See Appendix for details
}
\evaluation{As explained in the preceding section, the team size was reduced considerably, and only 1.0 scientist and a student assistant remain. The professor's position of 0.05 FTE was merely symbolic, as this was a requirement by the university administration when we presented the business model. 

There are additional costs which are not listed in the business model as the services would have been provided in-kind by the university. This would include hosting or legal services. These items were nevertheless established as calculatory costs (Hosting 120 EUR/month; legal fees 7000 EUR in the first year, 500 EUR in the years thereafter) However, they were never consolidated in the business model as such. The main reason for this is that the personnel costs present such an important part of the operations total budget that even 7000 EUR of legal costs do not really make a huge difference. Some estimates are given in the appendix. 
}
\othersolutions{
One could have (and probably should have) calculated the non-personnel costs in more detail.  In hindsight, however, we see that there were no legal costs in 2017 (besides the costs of registering the company, which are not born by the project). The design costs and IT costs, which were budgeted, were also not needed.  This is probably a sign of rather conservative estimations when we devised the model. 
}
  

\section{Investment}

\background{A business model requires you state the initial investment needed to start operations. For producing goods, this can be huge somes, for which financing has to be found. For IT-related projects, the sums are a lot less important these days.}
\langscisolution{
Only 6000 EUR are required for three work spaces (desk, chair, computers). 
}
\evaluation{Due to the reduction in staff, only two work spaces were required.}
\othersolutions{
One could have hired fully-equipped office space on a monthly basis, leading to less initial investment, but higher running costs. 
}
 

% \section{Pricing}
% 
% \background{1}
% \langscisolution{see appendix 1--6}
% \evaluation{3}
% \othersolutions{4}
%  
% 
% \section{Earnings forecast / profitability forecast}
% 
% \background{1}
% \langscisolution{see appendix 6}
% \evaluation{3}
% \othersolutions{4}
%  

\chapter{Outlook}

\background{The outlook section was mainly intended for the management of the Freie Universität to see the advantages. As such, it is not really part of the business model proper, but it is still nice to give a broader context of the ecology the project evolves in.}
\langscisolution{
Language Science Press strengthens the positioning of the Freie Universität as International Network University. It also strengthens the city of Berlin's reputation as an innovative place for doing science. The business model is fully specified and can easily be transferred to and be reused in other disciplines. A logical first step would be other fields of the humanities, e.g. archeology. The Freie Universität can acquire a first-mover position here and gain visibility far beyond the immediate region. The experiences gained regarding which approaches and methods do indeed work represent a major advantage for Berlin, the Freie Universität, and Language Science Press since there is no similar project elsewhere. 
}
\evaluation{Since the project was not realised within the Freie Universität, this section is moot. We are still working on the transferability, which is among other things what this very document is about.}
\othersolutions{
Instead of having a strategic outlook, one could also have focussed on financial matters and argue how the operations would generate profit, increase market share, help gain strategic advantages over the competition etc. This was not our intention. 
}
 
% \todo{square with spreadsheet}

\appendix
\chapter{Revenue streams}
The appendix gives the estimates for the various sources of revenue, and the expected expenditures.
\section{Print sales} 
\background{For every book sold in bookstores or online, a certain percentage remains with the middlemen, while another percentage goes to the publisher. The number of titles times the number of copies sold times the margin gives the total print revenue.}
\langscisolution{
  \fittable{\small
  \begin{tabular}{lrrrr}  
    \lsptoprule
	    &  books & sold copies per book&   margin per copy &  print revenue\\
    \midrule	
    Year 1  &   48               & 50                      &   10 €                                     &  24.000 €                       \\
    Year 2  &   54               & 50                      &   10 €                                     &  27.000 €                       \\
    Year 3  &   60               & 50                      &   10  €                                    &  30.000 €                       \\
    Year 4  &   66               & 50                      &   10 €                                     &  33.000 €                       \\
    Year 5  &   72               & 50                      &  10 €                                      & 36.000 €                        \\
    \lspbottomrule
    \end{tabular}
  }
}
\evaluation{In 2017, we published 26 books, not 48. We sold copies for a total of 4541 EUR + 206 GBP + 1228 USD, thus very short of the projected 24.000 EUR. The margin we calculated also did not take into account the free copies we have to provide to various libraries and archives, which cost us 3130 EUR}

\section{Book processing charges}  
\background{ 
In order to make content freely available to readers, one way is to charge authors so-called \textit{book processing charges} (BPCs). We projected that the BPCs would be 3.500 EUR (the cost price of the book), but that there would be waivers for authors who could not afford them. The amount of BPCs waived would diminish over time. The product of book output, BPC amount, and payment percentage gives a final projected revenue, ranging from 25.200 EUR to 88.200 EUR.
}
\newpage
\langscisolution{
  \fittable{\small
  \begin{tabular}{lrrrr}  
       \lsptoprule
         &   new books &  fees per book&  non-waived&  revenue\\
         \midrule
Year 1  &   48                 &    3.500 €                        &    15 \%                                                 &   25.200 €                             \\
Year 2  &   54                 &    3.500 €                        &    20 \%                                                 &   37.800 €                             \\
Year 3  &   60                 &    3.500 €                        &    25 \%                                                 &   52.500 €                             \\
Year 4  &   66                 &    3.500 €                        &    30 \%                                                 &   69.300 €                             \\
Year 5  &   72                 &    3.500 €                        &    35 \%                                                 &   88.200 €      \\
\lspbottomrule
    \end{tabular}
  }     
}
\evaluation{We do not really advertise BPCs. In 2017, we had one author who very much liked our service and paid 1000 EUR, and one project where we charged 1500 EUR for extra support, thus no BPCs proper. This amounts thus to 2/26, or roughly 8\%, with 92\% of all authors paying nothing at all (one could argue that this was the idea in the first place of course). The amount paid was also not 3500 EUR as projected,  but rather 1250 EUR. Altogether, this means that we made about 10\% of the expected revenue from BPCs.}
\newpage
\section{Institutional membership} 
\background{
Rather than charging readers or authors, one could also argue that the provision of a publication platform benefits all researchers, and that that effort should be covered jointly by institutions all over the world, since everybody profits. 
}
\langscisolution{
In the business model, we took over the model of \textit{Open Edition}, which has tiered membership prices, based on size of the institution and purchasing power of the country. 

%Open edition
{\small
\noindent
\begin{tabularx}{\textwidth}{lSSSSSSS}  
    \lsptoprule
students/uni &<500 &>500&>2.500&>10,000&>25,000&>50.000\\
staff/org.&<10&>10&>50&>250&>1,000&>5,000\\
\midrule
% Brutto-national-einkommen pro Kopf
Low         &100 €  &  250 €&  500 €&1,000 €&1,500 €&2.000 €\\
Lower Middle&500 €  &1,000 €&1,500 €&2,000 €&2,500 €&3,000 €\\
Upper Middle&1.500 €&2,000 €&2,500 €&3,000 €&3,500 €&4,000 €\\
High Income &2.500 €&3,000 €&3,500 €&4,000 €&4,500 €&5,000 €\\
\lspbottomrule
  \end{tabularx}
{\footnotesize 
  Low Income < 1,045 USD; Lower Middle Income < 4,125 USD; Upper Middle Income < 12,764 USD; High Income >12,764 USD;}
}

\medskip
We projected that the  number of supporting institutions would best be predicted by the number of series. Each series would bringe a new institional member. These members would contribute 3.500 EUR on average. As the number of members grows, the amount collected grows as well. 

\noindent
\begin{tabularx}{\textwidth}{Xrrrr}
\lsptoprule
        &  series&  members/series&       avg. contribution &     revenue\\
\midrule
Year 1 &  16               &       1                                    &        3.500 €                                               &     56.000 €                                     \\
Year 2 &  18               &       1                                    &        3.500 €                                               &     63.000 €                                     \\
Year 3 &  20               &       1                                    &        3.500 €                                               &     70.000 €                                     \\
Year 4 &  22               &       1                                    &        3.500 €                                               &     77.000 €                                     \\
Year 5 &  24               &       1                                    &        3.500 €                                               &     84.000 €                                     \\
\lspbottomrule
\end{tabularx}
}
\evaluation{
It turned out that the series-based estimations did not work out. Instead of growing slowly, we decided to go for 100 institutions from the beginning, at a much lower cost: 1000 EUR instead of 3500 EUR. The formula $100 \times 1000$ is much easier to communicate than the roster with size of institution and purchasing power of the country. The model $100 \times 1000$ means that we can collect 100.000 EUR right away, minus 15\% commission. 

Thus, instead of revenues of 56.000 EUR projected for the first year, we are at 0 EUR for the year 2016, and at 85.000 EUR for the year 2017. 
}

\newpage 
\section{Individual membership}  
\background{
Since readers and authors both profit, one should give them the opportunity to contribute on a regular basis as well. This should be done via membership in an association (\textit{Verein}).
}
\langscisolution{
The projection was that each new publication will triger a number of new supporters. Each supporter has a likelihood to become a paying member, paying a certain sum per year. We can thus calculate the total number of paying members in a given year and multiply this with the estimated average amount to arrive at the total revenue. 

\noindent
\begin{tabularx}{\textwidth}{Xrlllllr}
\lsptoprule 
          &  books   &  \rotatehead[2cm]{\mbox{\hspace*{-4mm}new members/book}}&  \rotatehead{\mbox{\hspace*{-4mm}new members}}&\rotatehead{\mbox{\hspace*{-4mm}conversion rate}}&\rotatehead{\mbox{\hspace*{-4mm}total members}}&\rotatehead{\mbox{\hspace*{-4mm}avg. contribution}}& revenue\\
\midrule
Year 1    &    48        &           5                  &                   600           &        .33                  &                      440               &            30 €                   &               13.200 €    \\
Year 2    &    54        &           5                  &                   100           &        .33                  &                      743               &            30 €                   &               22.300 €    \\
Year 3    &    60        &           5                  &                   100           &        .33                  &                      1.077             &            30 €                   &               32.300 €    \\
Year 4    &    66        &           5                  &                   100           &        .33                  &                      1.440             &            30 €                   &               43.200 €    \\
Year 5    &    72        &           5                  &                   100           &        .33                  &                      1.833             &            30 €                   &               55.000 €    \\
\lspbottomrule
\end{tabularx} 
}
\evaluation{
There is the association \textit{Open Science Press Support e.V}. with an annual membership fee of 20 EUR. This association is not actively promoted, and no new members have joined since its creation. 

Knowledge Unlatched had plans of having indidivual support pledges at the rate of 100 EUR/year in addition to institutional pledges. This turned out to be not very attractive to supporters. The LangSci team also did not really want to enter into financial dealings with their supporter base. Finally, 100 EUR/year is too low a sum for the overhead it generates in administration and acquisition, so this indidividual support pledge was no longer pursued and finally dropped. 7 indidividuals had pledged support. They were informed that they were welcome to make donations instead of their original pledges. This has the added advantage that there is no VAT on donations, and, furthermore, they are tax-deductible. 

The total sum raised by individual memberships is below 1000 EUR, and thus a far cry from the 13.000 EUR envisioned. 

}
\othersolutions{
One could have opted for a much more aggressive marketing campaign with flyers, mailings, referral schemes, special offers etc, coupled with a system tracking conversion rates and customer relations managements. This is too much overhead, and it is also not in line with the relations we want to have with our community. 
}

\newpage 
\section{Donations} 
\background{Donations are one-off payments. At the end of each pdf, there is a page reminding the authors that they have read the book for free, and that they could make a donation. Based on the number of books and an esimated number of donors with an estimated average amount, one can arrive at an estimated revenue from this channel.}
\langscisolution{
\begin{tabularx}{\textwidth}{Xrrrr}
\lsptoprule 
          &   books   & donors/book &   avg. amount     &  revenue\\
\midrule  
Year 1   &   48      &            20      &            10 €     &                      9.600 €  \\
Year 2   &   54      &            20      &            10 €     &                      10.800 € \\
Year 3   &   60      &            20      &            10 €     &                      12.000 € \\
Year 4   &   66      &            20      &            10 €     &                      13.200 € \\
Year 5   &   72      &            20      &            10 €     &                      14.400 € \\
\lspbottomrule
\end{tabularx}
}
\evaluation{
The number of books was 26, not 48, as already mentioned above. No one seemed to take note of the ad at the end of the book, at least we did not receive smaller donations of the sort we had envisaged at all. However, we did receive some larger donations, of more than 100 or 1000 EUR. 
}
\othersolutions{One could have been more nagging, with pop-ups on the website asking for donations and maybe even ads inside the book for non-members, similar to a freemium model. This would all have required programming effort, which would probably have been more expensive than the additional revenue which could have been gained from that. }

\chapter{Earnings forecast} 
\background{This is the complete calculation of all revenue and costs, based on estimates as to the number of series, the number of books, the number of download etc. Also confer the spreadsheet distributed together with this model.}
\langscisolution{
\fittable{
\begin{tabular}{ll rrrrr r}
\lsptoprule
 {} &  & \textbf{2016*} & \textbf{2017} & \textbf{2018} & \textbf{2019} & \textbf{2020} & \textbf{forw. proj.}\\
 \midrule
\multicolumn{8}{l}{ \textbf{Costs}}\\
\midrule
\textbf{1.} & \textbf{Personnel} &  146.070,00 € &  150.452,10 € &  170.062,27 € &  193.172,28 € &  217.515,83 € &  3,00\%\\
\textbf{2.} & \textbf{Travel} &  10.000,00 € &  10.300,00 € &  10.609,00 € &  10.927,27 € &  11.255,09 € &  3,00\%\\
\textbf{3.} & \textbf{Design and Marketing} &  3.000,00 € &  3.090,00 € &  3.182,70 € &  3.278,18 € &  3.376,53 € &  3,00\%\\
\textbf{4.} & \textbf{IT-Infrastructure} &  8.500,00 € &  8.755,00 € &  9.017,65 € &  9.288,18 € &  9.566,82 € &  3,00\%\\
\textbf{5.} & \textbf{Hosting} &  1.440,00 € &  1.483,20 € &  1.527,70 € &  1.573,53 € &  1.620,73 € &  3,00\%\\
\textbf{6.} & \textbf{Legal costs} &  7.000,00 € &  500,00 € &  500,00 € &  500,00 € &  500,00 € & ~\\
 \midrule 
\multicolumn{8}{l}{\textbf{Calculatory costs}}\\
\midrule
\textbf{7.} & \textbf{Office rent} &  9.000,00 € &  9.270,00 € &  9.548,10 € &  9.834,54 € &  10.129,58 € &  3,00\%\\
\textbf{8.} & \textbf{Office (Telco, stationary, postage} &  1.200,00 € &  1.236,00 € &  1.273,08 € &  1.311,27 € &  1.350,61 € &  3,00\%\\
\textbf{9.} & \textbf{Accounting} &  1.440,00 € &  1.483,20 € &  1.527,70 € &  1.573,53 € &  1.620,73 € &  3,00\%\\
\textbf{10.} & \textbf{Depreciation} &  1.230,77 € &  1.230,77 € &  1.230,77 € &  1.230,77 € &  1.230,77 € & ~\\
\textbf{11.} & \textbf{Insurance} &  500,00 € &  515,00 € &  530,45 € &  546,36 € &  562,75 € &  3,00\%\\
\textbf{12.} & \textbf{Risk  premium  (5\% of revenue)} &  6.400,00 € &  8.045,00 € &  9.840,00 € &  11.785,00 € &  13.880,00 € & ~\\
\midrule
\multicolumn{2}{c}{ \textbf{Total costs}} &  \textbf{195.780,77 €} &  \textbf{196.360,27 €} &  \textbf{218.849,41 €} &  \textbf{245.020,91 €} &  \textbf{272.609,45 €} & \textbf{~}\\
\end{tabular}
} 
{\tiny *  {Costs 2016}: The costs do not include the personnel costs still covered under the initial DFG funding (5/12 of yearly costs).
}
\newpage \noindent
\fittable{
\begin{tabular}{ll rrrrr r}
\lsptoprule
 {} &  & \textbf{2016*} & \textbf{2017} & \textbf{2018} & \textbf{2019} & \textbf{2020} \\
 \midrule
\multicolumn{5}{l}{ \textbf{Revenue}}\\
\midrule
\textbf{1.} & \textbf{Print sales} &  24.000,00 € &  27.000,00 € &  30.000,00 € &  33.000,00 € &  36.000,00 €\\
\textbf{2.} & \textbf{BPCs} &  25.200,00 € &  37.800,00 € &  52.500,00 € &  69.300,00 € &  88.200,00 €\\
\textbf{3.} & \textbf{Institutional supporters} &  56.000,00 € &  63.000,00 € &  70.000,00 € &  77.000,00 € &  84.000,00 €\\
\textbf{4.} & \textbf{Individual supporters} &  13.200,00 € &  22.300,00 € &  32.300,00 € &  43.200,00 € &  55.000,00 €\\
\textbf{5.} & \textbf{Donations} &  9.600,00 € &  10.800,00 € &  12.000,00 € &  13.200,00 € &  14.400,00 €\\
\midrule
\multicolumn{2}{c}{ \textbf{Total revenue}} &  \textbf{128.000,00 €} &  \textbf{160.900,00 €} &  \textbf{196.800,00 €} &  \textbf{235.700,00 €} &  \textbf{277.600,00 €}\\
\midrule\midrule
{} & \textbf{Net result} &  \textbf{{}-67.780,77 €} &  \textbf{{}-35.460,27 €} &  \textbf{{}-22.049,41 €} &  \textbf{{}-9.320,91 €} &  \textbf{4.990,55 €}\\
\lspbottomrule
\end{tabular}
}


{\scriptsize
\noindent
% \fittable{
\begin{tabularx}{\textwidth}{Qrrrrr}
\lsptoprule
{ \textbf{Estimates}} & \textbf{2016} & \textbf{2017} & \textbf{2018} & \textbf{2019} & \textbf{2020}\\
\midrule 
series &  \textbf{16} &  \textbf{18} &  \textbf{20} &  \textbf{22} &  \textbf{24}\\
new books per series&  \textbf{3} &  \textbf{3} &  \textbf{3} &  \textbf{3} &  \textbf{3}\\
total books &  \textbf{48} &  \textbf{54} &  \textbf{60} &  \textbf{66} &  \textbf{72}\\
print sales per new book&  50 &  50 &  50 &  50 &  50\\
print margin&  10,00 € &  10,00 € &  10,00 € &  10,00 € &  10,00 €\\
BPCs&  3.500,00 € &  3.500,00 € &  3.500,00 € &  3.500,00 € &  3.500,00 €\\
\% paid BPCs&  15,00\% &  20,00\% &  25,00\% &  30,00\% &  35,00\%\\
\# institutional members (1/series) &  16 &  18 &  20 &  22 &  24\\
avg. contribution IM&  3.500,00 € &  3.500,00 € &  3.500,00 € &  3.500,00 € &  3.500,00 €\\
supporters &  600 &  700 &  800 &  900 &  1000\\
individual supporters (5 per book, .33 new members per supporter) &  440 &  743 &  1.077 &  1.440 &  1.833\\
avg. contribution members &  30,00 € &  30,00 € &  30,00 € &  30,00 € &  30,00 €\\
donors &  960 &  1.080 &  1.200 &  1.320 &  1.440\\
avg. donation&  10,00 € &  10,00 € &  10,00 € &  10,00 € &  10,00 €\\
\lspbottomrule
\end{tabularx}
}

\bigskip 
\noindent
\fittable{
\begin{tabular}{lrrrrrrr}
\lsptoprule

\multicolumn{5}{l}{ \textbf{{Personnel planning 2016} - 2020}}\\
\midrule
& & \textbf{2016} & \textbf{2017} & \textbf{2018} & \textbf{2019} & \textbf{2020} & \textbf{forw. proj.}\\
\midrule 
\multicolumn{5}{l}{ Coordinator 1 (TV-L 13)}\\
\midrule
&  FTE &  100,00\% &  100,00\% &  100,00\% &  100,00\% &  100,00\% & ~\\
&  Salary &  65.920,00 € &  67.897,60 € &  69.934,53 € &  72.032,56 € &  74.193,54 € &  3,00\%\\
 &  Total &  65.920,00 € &  67.897,60 € &  69.934,53 € &  72.032,56 € &  74.193,54 € & ~\\
\\
\multicolumn{5}{l}{ Coordinator 2 (TV-L 13)}\\
\midrule
&  FTE &  50,00\% &  50,00\% &  50,00\% &  75,00\% &  100,00\% & ~\\
&  Salary &  65.920,00 € &  67.897,60 € &  69.934,53 € &  72.032,56 € &  74.193,54 € &  3,00\%\\
 &  Total &  32.960,00 € &  33.948,80 € &  34.967,26 € &  54.024,42 € &  74.193,54 € & ~\\
\\
\multicolumn{5}{l}{ Finances and fundraising(TV-L 13)}\\
\midrule
&  FTE &  50,00\% &  50,00\% &  50,00\% &  50,00\% &  50,00\% & ~\\
&  Salary &  65.920,00 € &  67.897,60 € &  69.934,53 € &  72.032,56 € &  74.193,54 € &  3,00\%\\
 &  Total &  32.960,00 € &  33.948,80 € &  34.967,26 € &  36.016,28 € &  37.096,77 € & ~\\
\\
\multicolumn{5}{l}{ Clerk (TV-L 9)}\\
\midrule
&  FTE &  25,00\% &  25,00\% &  50,00\% &  50,00\% &  50,00\% & ~\\
&  Salary &  56.920,00 € &  58.627,60 € &  60.386,43 € &  62.198,02 € &  64.063,96 € &  3,00\%\\
 &  Total &  14.230,00 € &  14.656,90 € &  30.193,21 € &  31.099,01 € &  32.031,98 € & ~\\
\\
\end{tabular}
}
\noindent
\fittable{
\begin{tabular}{lrrrrrrr}
\midrule
\multicolumn{5}{l}{ \textbf{{Personnel planning 2016} - 2020 ctd.}}\\
\midrule
& & \textbf{2016} & \textbf{2017} & \textbf{2018} & \textbf{2019} & \textbf{2020} & \textbf{forw. proj.}\\
\midrule 
\multicolumn{5}{l}{ Professor (W3)}\\
\midrule
&  FTE &  5,00\% &  5,00\% &  5,00\% &  5,00\% &  5,00\% & ~\\
&  Salary &  87.467,60 € &  90.091,63 € &  92.794,38 € &  95.578,21 € &  98.445,55 € &  3,00\%\\
 &  Total &  4.373,38 € &  4.504,58 € &  4.639,72 € &  4.778,91 € &  4.922,28 € & ~\\
\\
\multicolumn{5}{l}{ Student assistants (10h/week)}\\
\midrule
&  FTE &  50,00\% &  50,00\% &  50,00\% &  50,00\% &  50,00\% & ~\\
&  Salary &  6.500,00 € &  6.695,00 € &  6.895,85 € &  7.102,73 € &  7.315,81 € &  3,00\%\\
 &  Total &  3.250,00 € &  3.347,50 € &  3.447,93 € &  3.551,36 € &  3.657,90 € & ~\\
\midrule
\multicolumn{2}{l}{Grand total}  &  146.070,00 € &  150.452,10 € &  170.062,27 € &  193.172,28 € &  217.515,83 € & ~\\
\lspbottomrule
\end{tabular}
}

\bigskip
{\footnotesize
\noindent
\begin{tabularx}{\textwidth}{Xrrrrr}
\lsptoprule
Staff needed  & \textbf{2016} & \textbf{2017} & \textbf{2018} & \textbf{2019} & \textbf{2020}\\
\midrule
Scientist &  2,00 FTE &  2,00 FTE &  2,00 FTE &  2,25 FTE &  2,50 FTE\\
Clerk &  0,25 FTE &  0,25 FTE &  0,50 FTE &  0,50 FTE &  0,50 FTE\\
Professor &  0,05 FTE &  0,05 FTE &  0,05 FTE &  0,05 FTE &  0,05 FTE\\
Assistant &  0,50 FTE &  0,50 FTE &  0,50 FTE &  0,50 FTE &  0,50 FTE\\
\lspbottomrule
\end{tabularx} 
}
}
\evaluation{As detailed in the other sections of the appendix, most of the predictions were completely off the mark}

\addchap{Conclusion}
This document has detailed the business model for Language Science Press in 2015. The ideas, motivations and projection have been spelled out. They have then been squared with reality and evaluated with regard to their potential for a sustainable funding model. The main conclusion is that most projections were far off the mark, and that among the different revenue streams, institutional membership is by far the most interesting one with the other ones far behind. 

% \chapter{Series}
% \background{This is an overview of the number of series available at the time of creation of the model. Not relevant for business model in the strict sense, hence dropped}
% \langscisolution{2}
% \evaluation{3}
% \othersolutions{4}
\end{document}
